% Автоматаар үүсгэсэн — эх файл: docs/troubleshooting.md
\chapter{Түгээмэл алдаа ба шийдэл}

\hypertarget{docker-ux431ux430-ux441ux442ux435ux43a}{%
\section{Docker ба стек}\label{docker-ux431ux430-ux441ux442ux435ux43a}}

\begingroup\scriptsize\begin{longtable}[]{@{}
  >{\raggedright\arraybackslash}p{(\columnwidth - 4\tabcolsep) * \real{0.3333}}
  >{\raggedright\arraybackslash}p{(\columnwidth - 4\tabcolsep) * \real{0.3333}}
  >{\raggedright\arraybackslash}p{(\columnwidth - 4\tabcolsep) * \real{0.3333}}@{}}
\toprule
\begin{minipage}[b]{\linewidth}\raggedright
Алдаа
\end{minipage} & \begin{minipage}[b]{\linewidth}\raggedright
Шалтгаан
\end{minipage} & \begin{minipage}[b]{\linewidth}\raggedright
Шийдэл
\end{minipage} \\
\midrule
\endhead
\bottomrule
\endlastfoot
\texttt{no\ space\ left\ on\ device} & дүрс, лог, боть хуримтлагдсан & \texttt{docker\ system\ df} → \texttt{docker\ system\ prune\ -a} (боть хэвээр үлдэнэ) \\
\texttt{port\ is\ already\ allocated} & өөр процесс портыг эзэлсэн & \texttt{sudo\ ss\ -tlnp\ \textbackslash{}\textbar{}\ grep\ \textless{}порт\textgreater{}} \\
Контейнер \texttt{Restarting} давтана & OOM эсвэл тохиргооны алдаа & \texttt{docker\ inspect\ \textless{}нэр\textgreater{}\ \textbackslash{}\textbar{}\ grep\ -i\ oomkilled}; логийг үз \\
\texttt{exec\ format\ error} & дүрс arm64-д зориулагдаагүй & \texttt{docker\ manifest\ inspect\ \textless{}image\textgreater{}\ \textbackslash{}\textbar{}\ grep\ arm64} \\
\texttt{docker\ stats} санах ойг 0 гэнэ & cgroup memory идэвхгүй & SETUP Б.4 \\
Compose \texttt{.env} уншихгүй & \texttt{.env} файл \texttt{stack/} дотор биш & \texttt{docker\ compose\ config\ \textbackslash{}\textbar{}\ head} \\
\end{longtable}\endgroup

\hypertarget{thingsboard}{%
\section{ThingsBoard}\label{thingsboard}}

\begingroup\scriptsize\begin{longtable}[]{@{}
  >{\raggedright\arraybackslash}p{(\columnwidth - 2\tabcolsep) * \real{0.5000}}
  >{\raggedright\arraybackslash}p{(\columnwidth - 2\tabcolsep) * \real{0.5000}}@{}}
\toprule
\begin{minipage}[b]{\linewidth}\raggedright
Алдаа
\end{minipage} & \begin{minipage}[b]{\linewidth}\raggedright
Шийдэл
\end{minipage} \\
\midrule
\endhead
\bottomrule
\endlastfoot
3 минутаас удаан эхэлнэ & Pi дээр хэвийн. \texttt{start\_period:\ 180s} \\
\texttt{OOMKilled} & swap 2 GB эсэхийг шалга; \texttt{JAVA\_OPTS} дахь \texttt{-Xmx}-ийг 1200m болго \\
REST 401 & нууц үг солигдсон эсвэл токен хуучирсан \\
MQTT холбогдохгүй & ThingsBoard-ын MQTT нь \textbf{1884} порт (EMQX 1883) \\
Firmware төхөөрөмжид хүрэхгүй & Device → Manage firmware → багцыг оноох \\
\end{longtable}\endgroup

\hypertarget{emqx}{%
\section{EMQX}\label{emqx}}

\begingroup\scriptsize\begin{longtable}[]{@{}
  >{\raggedright\arraybackslash}p{(\columnwidth - 2\tabcolsep) * \real{0.5000}}
  >{\raggedright\arraybackslash}p{(\columnwidth - 2\tabcolsep) * \real{0.5000}}@{}}
\toprule
\begin{minipage}[b]{\linewidth}\raggedright
Алдаа
\end{minipage} & \begin{minipage}[b]{\linewidth}\raggedright
Шийдэл
\end{minipage} \\
\midrule
\endhead
\bottomrule
\endlastfoot
TLS \texttt{unknown\ ca} & \texttt{-\/-ca} буруу эсвэл сертификат өөр CA-гаас \\
TLS \texttt{hostname\ mismatch} & серверийн CN нь Pi-гийн нэртэй таарахгүй \\
Холболт 10000 дээр зогсоно & \texttt{EMQX\_\allowbreak{}LISTENERS\_\allowbreak{}\_\allowbreak{}TCP\_\allowbreak{}\_\allowbreak{}DEFAULT\_\allowbreak{}\_\allowbreak{}MAX\_\allowbreak{}CONNECTIONS} \\
Самбарт нэвтрэхгүй & \texttt{.env} дэх \texttt{EMQX\_\allowbreak{}DASHBOARD\_\allowbreak{}PASSWORD}; эхний эхлэлтийн дараа зөвхөн самбараас солино \\
\end{longtable}\endgroup

\hypertarget{influxdb-3-core}{%
\section{InfluxDB 3 Core}\label{influxdb-3-core}}

InfluxDB 3 Core-ийн CLI түргэн өөрчлөгддөг. Хэрэв \texttt{-\/-without-auth} эсвэл \texttt{-\/-object-store} таних тэмдэг ажиллахгүй бол:

\begin{Shaded}
\begin{Highlighting}[]
\ExtensionTok{docker}\NormalTok{ run }\AttributeTok{{-}{-}rm}\NormalTok{ influxdb:3.2{-}core influxdb3 serve }\AttributeTok{{-}{-}help}
\end{Highlighting}
\end{Shaded}

Шийдэгдэхгүй бол InfluxDB 2.7 руу шилжинэ:

\begin{Shaded}
\begin{Highlighting}[]
\ExtensionTok{docker}\NormalTok{ compose }\AttributeTok{{-}f}\NormalTok{ docker{-}compose.yml }\AttributeTok{{-}f}\NormalTok{ docker{-}compose.influx2.yml }\AttributeTok{{-}{-}profile}\NormalTok{ core up }\AttributeTok{{-}d}
\end{Highlighting}
\end{Shaded}

Энэ тохиолдолд Grafana-гийн өгөгдлийн эх сурвалж (SQL биш Flux), Node-RED-ийн бичих зангилаа өөр болно --- \texttt{lab05/reference/influx2/} дотор бэлэн хувилбар байгаа.

\hypertarget{ux441ux4afux43bux436ux44dux44d-ux431ux430-ux445ux43eux43bux431ux43eux43bux442}{%
\section{Сүлжээ ба холболт}\label{ux441ux4afux43bux436ux44dux44d-ux431ux430-ux445ux43eux43bux431ux43eux43bux442}}

\begingroup\scriptsize\begin{longtable}[]{@{}
  >{\raggedright\arraybackslash}p{(\columnwidth - 2\tabcolsep) * \real{0.5000}}
  >{\raggedright\arraybackslash}p{(\columnwidth - 2\tabcolsep) * \real{0.5000}}@{}}
\toprule
\begin{minipage}[b]{\linewidth}\raggedright
Алдаа
\end{minipage} & \begin{minipage}[b]{\linewidth}\raggedright
Шийдэл
\end{minipage} \\
\midrule
\endhead
\bottomrule
\endlastfoot
\texttt{pi-team03.local} олдохгүй & mDNS ажиллахгүй байна --- IP хаягаар хандана (\texttt{hostname\ -I}) \\
Хөтөч холбогдохгүй & VS Code PORTS таб эсвэл \texttt{ssh\ -L} туннел \\
MQTT \texttt{Connection\ refused} & Pi дээр брокер ажиллаж байна уу: \texttt{docker\ compose\ ps} \\
Гэнэт тасалдана & Wi-Fi эрчим хүч хэмнэх горим: \texttt{sudo\ iw\ wlan0\ set\ power\_save\ off} \\
\end{longtable}\endgroup

\hypertarget{ux445ux44dux43cux436ux438ux43bux442ux438ux439ux43d-ux447ux430ux43dux430ux440}{%
\section{Хэмжилтийн чанар}\label{ux445ux44dux43cux436ux438ux43bux442ux438ux439ux43d-ux447ux430ux43dux430ux440}}

\begin{itemize}
\tightlist
\item
  Хэмжилтийн өмнө \textbf{бусад ачааллыг зогсооно} (\texttt{htop}-оор шалга)
\item
  Гурваас доошгүй удаа давтаж дунджийг авна
\item
  \texttt{vcgencmd\ get\_throttled} нь \texttt{0x0} биш бол хэмжилтийг хүчингүй гэж үз
\item
  Цагийн синхрончлолгүй бол саатлын хэмжилт утгагүй (\texttt{timedatectl\ status})
\end{itemize}
